\section{Repair, pressure and transition errors on the generated range}

Write $R=\epsilon^{-1}$, $c=\nu^{-2}$,
$g(s)=\exp\{\gamma s/(1+\nu s)\}$ and
$N_\nu=e^{\lambda/\nu}\nu^{-\kappa}$, where $\Re\lambda=\gamma>0$.
The geometry satisfies
\[
 \epsilon=C_*\nu^r e^{-d_0/\nu},\qquad d_0=\gamma/3,
 \qquad \epsilon^3|N_\nu|\nu^{5+2\mu}\asymp1.
\]
All constants below are uniform for $\mu$ in a fixed compact subset
of $(0,1/4)$ and for the stipulated compact ranges of $C_*,r$.
The gain-band assertions refer to parameter triples satisfying the
displayed relation. At fixed $\nu,\epsilon$ the final spectral scan
is only an interval of width $O(|\log\nu|^{-1})$; no positive gain
bound is asserted on a nonshrinking spectral interval for that fixed
background.
$P_\nu$ denotes a fixed polynomial in $\nu^{-1}$. The exact
background is $U=V+r_b$, $V=U_{\rm raw}+E\alpha_{\rm bg}$, with the unprojected raw residual and
coefficient bounds of the background construction. In particular
\[
 \|r_b\|_{C_\tau L^{11/4}}+\|r_b\|_{W^{1,\infty}}
  +\sum_{j=1}^2\|D^jr_b\|_{Y_{11/4}}
 \le P_\nu\rho,\qquad
 \rho=\epsilon^{3/11}e^{m_0/\nu}.
\]
Fixed additional viscosity powers are included in $P_\nu$.
For an arbitrary axisymmetric physical velocity $w$ without swirl put
\[
 G(w)=R^{-1}\|g^{-1}w\|_{Y_2}.
\]

\subsection{Normalize the full old stress before inversion}

Let $\chi_o$ vanish on $s\le1.2/\nu$ and equal one on
$s\ge1.5/\nu$. The old source must be
\[
 T_o=\chi_o(U\otimes w+w\otimes U),
\]
including the repair. Define $W_\mu=-S_\mu T_o$,
$X=\mathcal C_\chi(W_\mu-W_{\mu+1})$,
$Y=L_{{\rm free},\mu+1}X$ and $h=[U,X]$, where
$\mathcal C_\chi$ is the divergence-free stream cutoff supported on
$s\le1/\nu$. All operators are applied in physical units here.
The tensor in the two shifted resolvents is identical.

The thin part is controlled by its first transverse moment. If
$w=\mathcal N_v b$ and $A=1.2/\nu$, weighted Cauchy--Schwarz gives
\[
 N_1(\chi_o M_vb)
 \le P_\nu A^{-1/2}\|b\|_H
 \le P_\nu |N_\nu|G(w).
\]
For the repair, use instead the global product estimate
\begin{equation}\label{repair:source}
 F_r:=\|\chi_o(r_b\otimes w+w\otimes r_b)\|_{L^2_\tau L^{22/19}}
 \le P_\nu\rho R|N_\nu|G(w).
\end{equation}
This follows from H\"older's inequality and $g\le e^{\gamma/\nu}$.
The fixed power relating $e^{\gamma/\nu}$ to $|N_\nu|$ is absorbed
in $P_\nu$. No moment is assigned to the repair source.

The general separated kernels give $X,DX,Y\in Y_2$ for both
source classes. For example the repair source-to-$L^2$ bound uses
the order-four kernel in $L^{11/7}$, with norm
$CD^{-23/11}$ outside distance $D$. Hence
$w\mapsto H_ow:=Y+h$ is bounded at fixed parameters.
Its output has young support. Consequently $H_o^2=0$ and
$(I+H_o)^{-1}=I-H_o$ by the disjoint old and young supports.

For $T=M_Uw$, the later full affine correction gives
$\widehat w=w+Y[T]+h[T]-h_{\mathcal LT}$. The added correction has
young support and depends only on old data. Hence
\begin{equation}\label{repair:old-invariance}
 \chi_oM_U\widehat w=\chi_oM_Uw.
\end{equation}
Thus every old-source estimate below applies literally to the new
unknown $w$. No gain norm of the possibly large young correction is
inserted into the old repair product.

\subsection{The repair source and exact pressure give a positive margin}

Let $\widetilde E(U;X,Y)$ be the full unprojected Ward residual,
\[
 2U\cdot\nabla Y-2\nu C_{\rm diff}
 +(U\otimes U):D^2X-(DX+DX^T)\nabla p.
\]
Here $p$ is the exact pressure. For the old repair source, the
pointwise gradient kernel has $L^{22/3}$ norm $CD^{-101/22}$.
At age $s\le1/\nu$ this gives
\[
 \|DX(s)\|_{L^2_\tau L^\infty_y}
 +\|DY(s)\|_{L^2_\tau L^\infty_y}
 \le P_\nu\epsilon^{101/22}s^{101/44}F_r,
\]
with the spatial supremum understood on that age surface. The
global gradient estimate into $L^{22/3}$ uses the kernel exponent
$11/3$ and is
\begin{equation}\label{repair:phase-gradient}
 \|DX\|_{L^\infty_\tau L^{22/3}}
 +\|DY\|_{L^\infty_\tau L^{22/3}}
 \le P_\nu\epsilon^{46/11}F_r.
\end{equation}
The phase supremum follows from the free equations for
$(\partial_\tau-y\cdot\nabla/2)DX$ and $DY$ and the normalized
phase estimate over $P\le1$. The instantaneous projected tensor
divergence is separated from the receiver and has the same kernel
bound. This argument differentiates no source in phase.

To estimate $V\cdot\nabla Y$, use $V=\sqrt{s}v$ and the physical
Jacobian $\pi R^2s^{-5/2}ds\,d\eta$. Its squared normalized norm
is bounded by
\[
 P_\nu\epsilon^{101/11}F_r^2
 \int_0^{1/\nu}s^{34/11}
             \|\sup_q|v(s,q,\cdot)|\|_2^2\,ds.
\]
The integral is finite at zero and polynomial in $\nu^{-1}$.
For the remaining thin-coefficient terms use
$R^{-1}(\|DV\|_2+\|D^2V\|_2+\|V\otimes V\|_2)\le P_\nu$.
Here their products with the returned derivatives use the ordinary
$L^\infty_\tau L^\infty_y$ bound on the cutoff fields, of size
$P_\nu\epsilon^{101/22}F_r$ for a first derivative. To obtain it,
apply the same phase norm argument at its $q=\infty$ endpoint:
dilation is an isometry on $L^\infty$, so along its characteristics
the norm difference is bounded by the integral of the $L^\infty$
norm of $BZ$. Young's inequality bounds both $Z$ and $BZ$ in
$L^2_\tau L^\infty_y$. Thus no product of two merely phase-$L^2$
factors is used for the global pressure term below.
Define the canonical pressure by
$\nabla p_V=-(I-\mathbb P)(V\cdot\nabla V)$. It obeys
$\|\nabla p_V\|_2\le C\|V\|_\infty\|DV\|_2$.
Therefore
\[
 R^{-1}\|\widetilde E(V;X,Y)\|_{Y_2}
 \le P_\nu\epsilon^{101/22}F_r
 \le P_\nu\epsilon^{85/22}|N_\nu|e^{m_0/\nu}G(w).
\]

The exact pressure correction is
\[
 \nabla(p-p_V)=-(I-\mathbb P)
 (V\cdot\nabla r_b+r_b\cdot\nabla V+r_b\cdot\nabla r_b),
\]
whose $Y_{11/4}$ norm is at most $P_\nu\rho$.
By H\"older's inequality, \eqref{repair:phase-gradient} bounds all
products with a phase-average coefficient, including
$D^2r_bDX$ and $DX\nabla(p-p_V)$. They contribute at most
\[
 P_\nu R^{-1}\rho\epsilon^{46/11}F_r
 \le P_\nu\epsilon^{52/11}|N_\nu|e^{2m_0/\nu}G(w).
\]
Terms with $D^2X$ have an additional curvature power. This proves
the old repair bound on the entire source interval, including every
large age and every spatial pressure tail.

For a thin old source, replacing the exact background pressure by
the flat lifted pressure has the same first error size. Indeed the
unprojected residual gives
\[
 \nabla(p_{\rm raw}-p_{\rm lift})
       =-(I-\mathbb P)F_{\rm raw}.
\]
The entrance convection terms and the explicit radial pressure
completion then give
$\|\nabla p-G_{\rm flat}\|_{Y_{11/4}}
 \le P_\nu\epsilon^{3/11}e^{m_0/\nu}$.
The general thin-source estimate is
$\|DX\|_{L^\infty_\tau L^{22/3}}
 \le P_\nu\epsilon^{57/22}N_1$.
H\"older's inequality and $R^{-1}=\epsilon$ give the factor
$\epsilon^{85/22}$. Nonpressure repair and entrance
coefficient changes are bounded in the same way or have a larger
curvature power. Thus the combined error $f_r$ satisfies
\begin{equation}\label{repair:body-error}
 R^{-1}\|f_r\|_{Y_2}
 \le P_\nu |N_\nu|G(w)
 \left(\epsilon^{85/22}e^{m_0/\nu}
       +\epsilon^{52/11}e^{2m_0/\nu}\right).
\end{equation}
All the unprojected forces here have young support, because every
term contains a derivative of $X$ or $Y$. Applying the Leray
projection does not change their curl, so these forces do not create
a direct vorticity-source moment at old ages.

\subsection{Both transition terms and the other age rectangles}

Let the near operator remove only the original far mask
$s<S$, $t>e^\delta s$, with $S=\epsilon^{-b}$ and
$0<\delta<\log(6/5)$. The young-old rectangle is contained in this
mask. Its complement is exactly the late-output part
$(1-\chi_y)F_{\rm far}$ and the early-source part
$\chi_yF_{\rm far}(1-\chi_o)$. Neither is identified with a local
causal inverse. The two differences produced by the stream cutoff
are also retained:
\[
 (\mathcal C_\chi-\chi)W_\mu,
 \qquad [L_{{\rm free},\mu+1},\mathcal C_\chi]
                   (W_\mu-W_{\mu+1}).
\]
They are supported on $0.8/\nu<s<1/\nu$.

On late output, $g(s)^{-1}\le e^{-4\gamma/(9\nu)}$; on early
source, $g(t)\le e^{3\gamma/(5\nu)}$. The general separated
kernel bounds, followed by the actual product estimate
\eqref{repair:source}, give the following weakest curvature powers
for the four source/receiver combinations, in the split stress norm:
\[
 \begin{array}{c|cc}
   &\text{thin receiver}&\text{scaled repair receiver}\\ \hline
 \text{thin source}&\epsilon^2&
      \epsilon^{49/22}e^{m_0/(2\nu)}\\
 \text{repair source}&\epsilon^{63/22}e^{m_0/\nu}&
      \epsilon^{79/22}e^{3m_0/(2\nu)}.
 \end{array}
\]
Each entry has a common fixed polynomial $P_\nu S^k$, and is
multiplied by $G(w)|N_\nu|e^{-4\gamma/(9\nu)}$ for late output,
or by $G(w)e^{3\gamma/(5\nu)}$ for early source. For example,
the thin-to-repair entry uses the order-four kernel from $L^2$ to
$L^{22/3}$, with off-diagonal norm $D^{-23/11}$, followed by
$\sqrt\rho$. The repair-to-thin entry uses the pointwise exponent
$79/22$, followed by $\rho R$. The repair-to-repair entry uses
the source $L^{22/19}$ norm and the $L^{11/3}$ kernel. The
outward stream primitive has the same norm powers, so both
transition terms have the late-output bound.
The thin-to-thin estimate at age zero does not use a first-moment
bound with a fixed positive lower source age. Use the mass estimate
$N_0(a)\le C(a^{-3/4}+a^{-1/2})\|T\|_{Z_1}$ proved below instead.
The separated pointwise velocity is at most
$P_\nu\epsilon^2s^2N_0(e^\delta s)$. After the normalized velocity
lift and thin multiplication, the square root of the output stress
age density is bounded by
\[
 P_\nu\epsilon^2s^{5/4}V_m(s)N_0(e^\delta s)
 \le P_\nu\epsilon^2(s^{1/2}+s^{3/4})V_m(s)\|T\|_{Z_1}.
\]
Its square is integrable at zero, and the integral up to $S$ is
polynomial in $S$. This supplies the first table entry on the full
age interval.

The two weakest exponential margins are
\[
 2d_0-\gamma+4\gamma/9=\gamma/9,
 \qquad 2d_0-3\gamma/5=\gamma/15.
\]
For completeness the thin-to-repair margins are
$37\gamma/198$ and $47\gamma/330$, before its $m_0/2$ loss;
the repair-to-thin margins are $79\gamma/198$ and
$39\gamma/110$, before its $m_0$ loss. Thus, if $\delta T_a$
is the total stress correction from these regions and both stream
transitions,
\begin{equation}\label{repair:age-error}
 \|g^{-1}\delta T_a\|_Z
 \le P_\nu
  \exp\!\left[-\frac{\gamma/15-kbd_0-3m_0/2}{\nu}\right]G(w).
\end{equation}

\subsection{The whole thin geometric residual and the actual affine lift}

We now include the terms remaining in the exact transformed equation,
not just the repair and age errors. Write $T=(T_1,T_2)=M_Uw$ in
the split stress space, so $T_1=M_Vw$ is the thin component and the
physical repair tensor represented by $T_2$ is $M_{r_b}w$.
Let $X_1,Y_1$ denote the response to $\chi_oT_1$ alone and let
$X^{\rm mt},Y^{\rm mt}=D_{\mu+1}X^{\rm mt}$ be its phase-mean,
zero-transverse-width Newton response, with the same stream cutoff.
The uncut Newton fields are harmonic outside the old source disk;
their shifted dilation identity therefore gives the same displayed
commutator formula for $Y^{\rm mt}$. In particular the Laplacian
of the stream-cutoff field is retained even though the uncut fields
are harmonic there.
Set
\[
 A_T(s)=\mathcal L T_1(s)
       =D_yX^{\rm mt}(R/\sqrt s,0,0),\qquad
 B_T=A_T'+\mu A_T/s.
\]
The full Cartesian matrix is retained, including its hoop entry.
The first stream cutoff $\chi_y$ is supported in $s\le1/\nu$;
the affine source cutoff $\chi_f$ equals one there and vanishes for
$s\ge1.1/\nu$. Hence $\chi_f A_T=A_T$ exactly. In particular the
source $F_{A_T}$ uses $B_T$, without a second cutoff derivative.

Put $N=N_1(\chi_oT_1)$. The explicit separated kernels imply,
for $0\le j\le6$,
\begin{equation}\label{repair:actual-jet}
 |(s\partial_s)^j A_T(s)|\le C_j\epsilon^3s^{5/2}N,
 \qquad A_T(s)=0\quad(s\ge1/\nu).
\end{equation}
Indeed the gradient Newton kernel has degree $-5$, and the ratio of
the old source radius to the young receiver radius is at most
$1/\sqrt{1.2}<1$. Each $s\partial_s=-r\partial_r/2$ derivative
keeps degree $-5$ after its factor $r$ is included. The outward
stream primitive has the same degree. Fixed conormal derivatives
of $\widehat\chi_y(\nu s)$ are bounded. This proves the displayed
bound directly, without first passing through the exponentially
weighted smooth-source norm.

For clarity we give the geometric comparison that is needed for
arbitrary thin inputs. One phase primitive in the dilation variable,
with its dilation-one boundary term retained, gives
\[
 \|D_yX_1-D_yX^{\rm mt}\|_{L^2_\tau}
 +\|D_yY_1-D_yY^{\rm mt}\|_{L^2_\tau}
       \le P_\nu\epsilon^4s^{5/2}N.
\]
The same bound includes the finite transverse width and the
heat-to-Newton error: Taylor's formula costs the first transverse
moment, and the strict radial gap holds along every Taylor segment.
The initially stronger age power $s^3$ is bounded by
$P_\nu s^{5/2}$ on the young interval. For $Y$ use the exact formula
\[
 Y=\mathcal C_{\chi_y}W_\mu+
 [D_{\mu+1},\mathcal C_{\chi_y}](W_\mu-W_{\mu+1}).
\]
Apply the spatial kernel comparison to these terms separately.
Thus no phase derivative of the improved bound, or of the source,
is taken. These are normalized phase $L^2$ estimates; the raw-train
coefficients multiplying them are bounded in phase with all the
finite transverse weights used below.

The mean transverse Taylor estimate and the Hessian estimate are
\[
 \begin{split}
 |D_yX^{\rm mt}(s,\eta)-A_T(s)|
 +|D_yY^{\rm mt}(s,\eta)-sB_j(s)|
   &\le P_\nu\epsilon^4s^{5/2}|\eta|N,\\
 \|D_y^2X_1(s,\eta)\|_{L^2_\tau}
   &\le P_\nu\epsilon^4s^3N,
 \qquad B_j=s^{-1}D_yY^{\rm mt}(R/\sqrt s,0,0).
 \end{split}
\]
There is an exact clock correction, including on the cutoff annulus:
\begin{equation}\label{repair:exact-jet-clock}
 B_T=B_j+\frac\nu sD_y\Delta_yX^{\rm mt}(R/\sqrt s,0,0),
 \qquad |B_T-B_j|\le P_\nu\nu\epsilon^5s^{5/2}N.
\end{equation}
To check its sign and constant, differentiate
$Y^{\rm mt}=(B-\nu\Delta+\mu+1/2)X^{\rm mt}$, use
$D_yB=BD_y-D_y/2$, and evaluate the phase-independent matrix at
$z=0$, where $B D_yX^{\rm mt}=sA_T'$. All stream-cutoff terms
are already in $Y^{\rm mt}$; none are omitted from this identity.

The exact ring-chart rules give
$U_{\rm raw}=\sqrt s(v+\epsilon H e_z)$,
$DU_{\rm raw}=s(D_{\rm flat}v+\epsilon C_1)$ and
$D^2U_{\rm raw}=s^{3/2}(D_{\rm flat}^2v+\epsilon C_2)$.
Here $H,C_1,C_2$ have the finite weighted train bounds supplied by
the explicit chart differentiation and the young train construction.
Substitution into the four terms of the Ward residual gives the
affine expression
\[
 s^{3/2}\{2B_Tv+2\nu A_T:D_{\rm flat}^2v
                      -(A_T+A_T^T)\nabla_{\rm flat}p_{\rm cyl}\}.
\]
The error in $2U\cdot\nabla Y$ has size
$P_\nu\epsilon^4s^3N$ times a weighted train coefficient.
The errors in $DXD^2U$, $DX\nabla p$, $DUD^2X$ and
$(U\otimes U):D^2X$ have size
$P_\nu\epsilon^4s^4N$. The clock error has the smaller size
$P_\nu\nu\epsilon^5s^4N$. In physical units the squared body
norm divided by $R^2$ uses the density $\pi s^{-5/2}ds\,d\eta$;
the remaining powers $s^{7/2}$ and $s^{11/2}$ are integrable at
zero and give only viscosity polynomials at $1/\nu$. This also
covers coordinate factors $q-jc$ in an alternative fixed-center
Taylor expansion, since they are Gaussian moments of fixed order.

Replacing the actual train and its decaying pressure by the summed
Gaussian profiles costs $P_\nu\epsilon^3\delta_{\rm base}N$.
Here the young-train comparison supplies the quantitative envelope
\[
 \delta_{\rm base}\le
       P_\nu\exp(C_G/\nu-b_G/\nu^4),\qquad b_G>0,
\]
for the fixed derivative and transverse-weight budget on
$0\le s\le1.1/\nu$, with $C_G,b_G$ independent of the small
parameters. It follows directly from the Gaussian train separation
$c=\nu^{-2}$. This is the bound used when absorbing its contribution
in an exponential curvature margin below.
The exact-pressure and repair-coefficient errors were estimated
above, using the ordinary phase-supremum gradient bound, not a
phase derivative of the extra-$\epsilon$ comparison. Consequently,
with $E_{A_T}$ the Gaussian residual of the physical affine identity,
\begin{equation}\label{repair:thin-complete}
 R^{-1}\|E(U;X_1,Y_1)-E_{A_T}\|_{Y_2}
 \le P_\nu N\bigl(\epsilon^4+
       \epsilon^{85/22}e^{m_0/\nu}
                         +\epsilon^3\delta_{\rm base}\bigr).
\end{equation}
The $85/22$ term retains the entire Cartesian pressure correction.

Apply the explicit physical affine lifting identity to the actual
jet $A_T$. Its proof uses at most three conormal derivatives of the
matrix, so \eqref{repair:actual-jet} gives directly
\begin{equation}\label{repair:actual-affine-remainder}
 R^{-1}\|r_{A_T}\|_{Y_2}
 \le P_\nu\epsilon^3N
       (\epsilon+\rho+\delta_{\rm base}).
\end{equation}
This follows by inserting the displayed conormal bounds into each
Gaussian product in that identity. Geometric terms have one extra
$\epsilon$, the exact-background convection difference is
$(U-U_{\rm raw})\cdot\nabla\widehat h_{A_T}
+\widehat h_{A_T}\cdot\nabla(U-U_{\rm raw})$, and its bound is
$P_\nu\rho$ times the normalized $H^1$ bracket bound. The explicit
lifted pressure potential is a global gradient plus its
$\epsilon$ radial completion. Thus no projection tail is dropped.
There is no factor $e^{1.1a_s/\nu}$ in this actual-range estimate;
that factor is needed only for a generic source in $\mathscr A$.

Since $N\le P_\nu|N_\nu|\|g^{-1}T\|_Z$, the new thin and
affine body errors have exponential margins at least
$\gamma/3$, $19\gamma/66-m_0$ and $\gamma/11-m_0$.
All exceed the weakest age stress margin $\gamma/15$ when the
stipulated losses are small.

\subsection{The direct near solver and an actual-range pair norm}

The direct zero-entrance near inverse satisfies
\[
 \|g^{-1}D_L^{-1}g\|_{Z\to Z}\le P_\nu,
 \quad
 \|g^{-1}M_UR_\mu f\|_Z\le P_\nu R^{-1}\|f\|_{Y_2},
 \quad
 \|g^{-1}S_{\rm near}g\|_{H_T\to H}
       \le P_\nu e^{\eta_\delta/\nu}.
\]
For $D_LT=M_UR_\mu f$, the exact equation gives
\[
 T=M_Uw,\qquad w=R_\mu f-S_{\rm near}T.
\]
It follows that $G(w)\le P_\nu e^{\eta_\delta/\nu}
R^{-1}\|f\|_{Y_2}$. No homogeneous entrance appears. For a general
product datum $M_Uv$, however, the reconstruction is
$w=v-S_{\rm near}T$. Its velocity norm must retain the norm of $v$;
a small stress norm does not imply a small velocity witness.

On the distinguished gain band, \eqref{repair:body-error} has
operator size at most
\[
 P_\nu\left[e^{-(19\gamma/66-m_0)/\nu}
             +e^{-(19\gamma/33-2m_0)/\nu}\right]
\]
from the gain-energy norm to the normalized body norm. For the full
actual affine bracket, \eqref{repair:actual-affine-remainder} and
$\epsilon^3N\le P_\nu A$ give the background-error margin
$\gamma/11-m_0$. Geometric lift errors with an extra $\epsilon$
have the larger margin $\gamma/3$.

Here is the additional estimate needed to propagate actual-range
witnesses. Let $A=\|g^{-1}T\|_Z$, with $T=M_Uw$. For an unweighted
thin stress and $a>0$, its old physical mass divided by $\pi R^2$ is
\[
 N_0(a)=\int_a^\infty t^{-3/2}\int\|T(t,\cdot,\eta)\|_{L^2_q}
                  \,d\eta\,dt
 \le C(a^{-3/4}+a^{-1/2})\|T\|_{Z_1}.
\]
Indeed transverse Cauchy--Schwarz contributes $(1+t)^{1/4}$,
and the remaining squared age integral is
$\int_a^\infty t^{-5/2}(1+t)^{1/2}dt$. Partition the output ages
into logarithmic bins of width $h<\delta/2$. On a bin of upper age
$s_j$, all source ages exceed $e^{\delta-h}s_j$, so there is a
uniform radial gap $D_j\asymp R/\sqrt{s_j}$. The order-four
kernel has spatial $L^2$ norm $CD_j^{-5/2}$. Minkowski and the
mass estimate therefore bound the normalized velocity on that bin by
\[
 C\epsilon^{3/2}(s_j^{1/2}+s_j^{3/4})\|T\|_{Z_1}.
\]
Summing geometric bins down to age zero is convergent. Up to a fixed
polynomial in $S$, the far thin-source velocity operator has norm
$P_\nu S^k\epsilon^{3/2}$ from $Z_1$ to the normalized physical
$L^2$ space. The outward stream primitive obeys the same bound.
This proves the estimate for both stream transitions as well.

Thus the late-output thin velocity has possible growth
$e^{\gamma/(18\nu)}$ after gain weighting, and the early-source
thin velocity has possible growth $e^{\gamma/(10\nu)}$:
\[
 \tfrac32d_0-\gamma+4\gamma/9=-\gamma/18,
 \qquad \tfrac32d_0-3\gamma/5=-\gamma/10.
\]
For an actual repair source, Young's $L^{22/19}\to L^2$ bound
instead gives normalized velocity
$P_\nu S^k\rho\epsilon^{23/11}$ times its gain. Its late and early
margins are respectively $23\gamma/99$ and $31\gamma/165$, before
the $m_0$ loss. These bounds include the full old interval.

Put $\omega_b=kbd_0$, increasing the fixed $k$ if necessary.
The error families proved in this section have a physical velocity
witness $v_e$, with stress $M_Uv_e$, satisfying
\begin{align*}
 \|g^{-1}M_Uv_e\|_Z&\le\delta_A A+\delta_G G(w),\\
 R^{-1}\|g^{-1}v_e\|_{Y_2}&\le\chi_A A+\chi_G G(w),
\end{align*}
where it suffices to take
\begin{align*}
 \delta_A&=P_\nu e^{-(\gamma/15-\omega_b-3m_0/2)/\nu},&
 \delta_G&=P_\nu e^{-(19\gamma/66-\omega_b-3m_0/2)/\nu},\\
 \chi_A&=P_\nu e^{(\gamma/10+\omega_b+3m_0/2)/\nu},&
 \chi_G&=P_\nu e^{-(31\gamma/165-\omega_b-3m_0/2)/\nu}.
\end{align*}
The pressure and Ward body-force contributions use $v_e=R_\mu f_e$,
whose energy norm is bounded by the normalized body norm. Their thin
source bounds use $A$, whereas the repair source bounds use $G(w)$.
The full affine bracket's background error is bounded by
\eqref{repair:actual-affine-remainder} and is included in $\delta_A$.
The actual jet has the direct bound \eqref{repair:actual-jet}; on
the unit gain band $\epsilon^3N_1\le P_\nu A$.

Let $Q=P_\nu$ and $K_s=P_\nu e^{\eta_\delta/\nu}$
be the direct near bounds. For one correction solve
$D_LT_+=-M_Uv_e$ and set $w_+=-v_e-S_{\rm near}T_+$.
Then $T_+=M_Uw_+$ exactly, and
\[
 A_+\le Q(\delta_A A+\delta_G G),\qquad
 G_+\le\chi_A A+\chi_G G+K_s A_+.
\]
This proves closure on the generated range, with no replacement of
the velocity witness by its stress norm. More precisely, use the
closed graph space
\begin{equation}\label{repair:pair-space}
 \mathcal E_\nu=\{(T,w):g^{-1}T\in Z,\ g^{-1}w\in Y_2,
                       \ T=M_Uw\}.
\end{equation}
The velocities in this graph have the required orthogonal symmetry
but need not be divergence-free. Boundedness of the thin multiplier and
the repair multiplier proves that the graph is closed. Thus the pair
norm below makes it a Banach space. The displayed update preserves
this graph exactly; the divergence-free condition is recovered from the full
physical kernel equation, not from the masked near reconstruction.
Choose $a=\gamma/8$ and use the pair norm $A+e^{-a/\nu}G$.
The four matrix coefficients are
\[
 Q\delta_A,\quad Q\delta_Ge^{a/\nu},\quad
 e^{-a/\nu}\chi_A+e^{-a/\nu}K_sQ\delta_A,\quad
 \chi_G+K_sQ\delta_G.
\]
All tend to zero exponentially if
$\eta_\delta,\omega_b,m_0$ are sufficiently small. For example
each less than $\gamma/1000$ leaves a margin larger than
$\gamma/50$ before absorbing fixed viscosity powers. The narrowest
new margin is $a-\gamma/10=\gamma/40$. An initial body-forced
column has pair norm at most $P_\nu e^{\eta_\delta/\nu}$, so
the resulting correction series also tends to zero under the bounded
old smooth-jet observation.

\begin{proposition}\label{repair:complete-graph-error}
Let $\mathcal E_\nu$ have the pair norm just defined. Define the exact error
velocity and its graph correction by
\[
 v_Q(T)=S_{\rm far}T+Y[T]
       +R_\mu\mathbb P\{E[T]-E_{\mathcal LT}-r_{\mathcal LT}\},
 \qquad \mathcal Q(T,w)=-\mathcal J(v_Q(T)).
\]
Here $E[T]$ uses the full old source, $\mathcal LT$ uses its thin
part, and all projections and cutoffs are the ones specified above.
Every term of this exact $v_Q$ is covered by the preceding estimates.
If $\eta_\delta,\omega_b,m_0<\gamma/1000$ and the background
and train envelopes satisfy the stated bounds, then, for sufficiently
small $\nu$ on the distinguished gain band,
\[
 \|\mathcal Q\|_{\mathcal E_\nu\to\mathcal E_\nu}
           \le P_\nu e^{-\gamma/(100\nu)}.
\]
For each fixed geometry this is a real operator-norm-continuous
family in $\mu$ on the stipulated compact interval.
\end{proposition}
\begin{proof}
The decomposition of $S_{\rm far}T+Y[T]$ consists exactly of the
early-source and late-output regions and the two stream transitions.
For the body term, linearity splits the full old response into its
thin and repair responses. The repair response is
\eqref{repair:body-error}; its thin counterpart is
\eqref{repair:thin-complete}; the final subtraction is
\eqref{repair:actual-affine-remainder}. These exhaust the displayed
definition, including finite width, oscillatory phase, the exact
matrix clock, pressure and the physical affine lift. The graph
lift gives precisely the four pair-norm coefficients above.
Their smallest unperturbed exponent is $\gamma/40$. The specified
losses leave more than $\gamma/50$, hence the stated weaker
$\gamma/100$ bound absorbs every fixed viscosity polynomial.
Continuity follows from the resolvent identity, the dominated
separated dilation kernels and the polynomial near inverse; the
products and finite conormal affine construction are continuous.
No differentiation of a rough source in phase is required.
\end{proof}
